% JEFE 2 — Tablas infernales: tabularray, celdas combinadas en filas y columnas, colores alternados,
% columnas decimales, diagonales, tablas largas en hoja horizontal, tablas generadas desde datos
% y tablas dentro de cajas. Todo lo que suele romperse al pasar de Overleaf a otro editor.
\documentclass[11pt,a4paper]{article}
\usepackage[T1]{fontenc}
\usepackage[spanish,es-noshorthands,es-tabla]{babel}
\usepackage{lmodern}
\usepackage[margin=2cm]{geometry}
\usepackage{amsmath}
\usepackage[table]{xcolor}
\usepackage{array,booktabs,multirow,makecell,tabularx,longtable,hhline,diagbox,threeparttable}
\usepackage{tabularray}
\UseTblrLibrary{booktabs,siunitx}
\usepackage{siunitx}
\sisetup{output-decimal-marker={,}}
\usepackage{pdflscape,rotating}
\usepackage{pgfplotstable}
\pgfplotsset{compat=1.18}
\usepackage[most]{tcolorbox}
\usepackage{caption}
\newcolumntype{C}[1]{>{\centering\arraybackslash}m{#1}}
\newcolumntype{Y}{>{\raggedleft\arraybackslash}X}

\begin{document}
\section*{1. Celdas combinadas en las dos direcciones}
\begin{table}[h]
\centering
\caption{\texttt{multirow} + \texttt{multicolumn} + colores + \texttt{hhline}.}
\begin{tabular}{|c|c|c|c|c|}
  \hhline{|-|-|-|-|-|}
  \rowcolor{blue!15}
  \multirow{2}{*}{\diagbox{Material}{Ensayo}} & \multicolumn{2}{c|}{Tracción} & \multicolumn{2}{c|}{Compresión}\\
  \hhline{|~|-|-|-|-|}
  \rowcolor{blue!15}
   & Máx. & Mín. & Máx. & Mín.\\
  \hhline{|=|=|=|=|=|}
  Acero & 400 & 250 & \multicolumn{2}{c|}{\cellcolor{red!20}no aplica}\\ \hline
  \multirow{3}{*}{\makecell{Concreto\\(3 mezclas)}} & 3,1 & 2,0 & 35 & 28\\ \cline{2-5}
   & 3,4 & 2,2 & \multirow{2}{*}{\cellcolor{green!15}38} & 30\\ \cline{2-3}\cline{5-5}
   & 3,0 & 1,9 & & 27\\ \hline
\end{tabular}
\end{table}

\section*{2. tabularray con todo}
\begin{table}[h]
\centering
\caption{\texttt{tblr}: combinadas con \texttt{SetCell}, columnas \texttt{si}, filas alternadas.}
\begin{tblr}{
  colspec = {l Q[si={table-format=3.2},c] Q[si={table-format=2.1},c] X[c]},
  row{1} = {font=\bfseries, bg=teal!20},
  row{even[2]} = {bg=gray!8},
  hline{1,Z} = {1pt}, hline{2} = {0.6pt},
  cell{2}{4} = {r=3}{c, bg=yellow!20},
  cell{6}{1} = {c=2}{l},
  width = 0.9\linewidth,
}
Componente & {{{Precio (S/)}}} & {{{Cantidad}}} & Observación\\
Resistencia & 0.50 & 20.0 & Compradas juntas en la librería, mismo proveedor y mismo día\\
Capacitor & 1.20 & 10.0 & \\
Inductor & 12.75 & 2.0 & \\
Placa & 25.00 & 1.0 & Perforada\\
Total (con IGV incluido) & & 3.5 & {Texto largo que se\\parte en dos líneas}\\
\end{tblr}
\end{table}

\section*{3. Tabla generada desde datos (pgfplotstable)}
\pgfplotstableread[col sep=comma]{
n,x,cuadrado,raiz
1,0.5,0.25,0.7071
2,1.0,1.00,1.0000
3,1.5,2.25,1.2247
4,2.0,4.00,1.4142
5,2.5,6.25,1.5811
6,3.0,9.00,1.7321
}\datos
\begin{center}
\pgfplotstabletypeset[
  columns={n,x,cuadrado,raiz},
  columns/n/.style={column name={$n$}, int detect},
  columns/x/.style={column name={$x$}, fixed, precision=1, zerofill},
  columns/cuadrado/.style={column name={$x^2$}, fixed, precision=2, zerofill},
  columns/raiz/.style={column name={$\sqrt{x}$}, fixed, precision=4, zerofill},
  create on use/suma/.style={create col/expr={\thisrow{x}+\thisrow{cuadrado}}},
  columns/suma/.style={column name={$x+x^2$}, fixed, precision=2},
  columns={n,x,cuadrado,raiz,suma},
  every head row/.style={before row=\toprule, after row=\midrule},
  every last row/.style={after row=\bottomrule},
  every even row/.style={before row={\rowcolor{orange!10}}},
  use comma, 1000 sep={\,}
]\datos
\end{center}

\section*{4. Tabla dentro de una caja que se parte entre páginas}
\begin{tcolorbox}[breakable, enhanced, colback=blue!3, colframe=blue!50!black,
  title={Caja con tabla \texttt{tabularx} y fórmulas}]
\begin{tabularx}{\linewidth}{>{\bfseries}l X Y}
  \toprule
  Ley & Enunciado & Fórmula\\
  \midrule
  Ohm & La corriente es proporcional al voltaje. & $V = IR$\\
  Kirchhoff (corrientes) & La suma de corrientes en un nodo es cero. & $\sum_k i_k = 0$\\
  Kirchhoff (voltajes) & La suma de voltajes en una malla es cero. & $\sum_k v_k = 0$\\
  Faraday & El voltaje inducido es la derivada del flujo, con signo menos. & $\varepsilon = -\dfrac{d\Phi}{dt}$\\
  \bottomrule
\end{tabularx}
\end{tcolorbox}

\section*{5. Tablas anidadas y celdas con varias líneas}
\begin{center}
\begin{tabular}{|C{3cm}|C{5cm}|}
  \hline
  \thead{Exterior\\(encabezado)} & \thead{Tabla interna}\\ \hline
  Celda simple &
  \begin{tabular}{@{}cc@{}} a & b\\ c & \begin{tabular}{|c|}\hline d\\ \hline\end{tabular}\end{tabular}\\ \hline
  \makecell{Una\\dos\\tres} & $\begin{pmatrix} 1 & 2\\ 3 & 4 \end{pmatrix}$ \\ \hline
\end{tabular}
\end{center}

\begin{landscape}
\section*{6. Tabla larga en hoja horizontal}
\begin{longtable}{>{\bfseries}c *{8}{S[table-format=3.1]}}
  \caption{Temperaturas (\si{\celsius}) — 60 filas en hoja horizontal.}\\
  \toprule
  Día & {Lun} & {Mar} & {Mié} & {Jue} & {Vie} & {Sáb} & {Dom} & {Prom.}\\
  \midrule\endfirsthead
  \toprule
  Día & {Lun} & {Mar} & {Mié} & {Jue} & {Vie} & {Sáb} & {Dom} & {Prom.}\\
  \midrule\endhead
  \bottomrule\endlastfoot
  1 & 19.2 & 21.3 & 23.4 & 22.3 & 24.0 & 25.6 & 23.9 & 22.8\\
  2 & 20.3 & 22.4 & 24.3 & 23.0 & 24.6 & 25.9 & 24.0 & 23.5\\
  3 & 21.4 & 23.3 & 25.0 & 23.6 & 24.9 & 26.0 & 23.8 & 24.0\\
  4 & 22.3 & 24.0 & 25.6 & 23.9 & 25.0 & 25.8 & 23.5 & 24.3\\
  5 & 23.0 & 24.6 & 25.9 & 24.0 & 24.8 & 25.5 & 22.9 & 24.4\\
  6 & 23.6 & 24.9 & 26.0 & 23.8 & 24.5 & 24.9 & 22.1 & 24.2\\
  7 & 23.9 & 25.0 & 25.8 & 23.5 & 23.9 & 24.1 & 21.1 & 23.9\\
  8 & 24.0 & 24.8 & 25.5 & 22.9 & 23.1 & 23.1 & 20.0 & 23.3\\
  9 & 23.8 & 24.5 & 24.9 & 22.1 & 22.1 & 22.0 & 18.8 & 22.6\\
  10 & 23.5 & 23.9 & 24.1 & 21.1 & 21.0 & 20.8 & 17.6 & 21.7\\ \midrule
  11 & 22.9 & 23.1 & 23.1 & 20.0 & 19.8 & 19.6 & 16.5 & 20.7\\
  12 & 22.1 & 22.1 & 22.0 & 18.8 & 18.6 & 18.5 & 15.3 & 19.6\\
  13 & 21.1 & 21.0 & 20.8 & 17.6 & 17.5 & 17.3 & 14.3 & 18.5\\
  14 & 20.0 & 19.8 & 19.6 & 16.5 & 16.3 & 16.3 & 13.5 & 17.4\\
  15 & 18.8 & 18.6 & 18.5 & 15.3 & 15.3 & 15.5 & 12.8 & 16.4\\
  16 & 17.6 & 17.5 & 17.3 & 14.3 & 14.5 & 14.8 & 12.3 & 15.5\\
  17 & 16.5 & 16.3 & 16.3 & 13.5 & 13.8 & 14.3 & 12.0 & 14.7\\
  18 & 15.3 & 15.3 & 15.5 & 12.8 & 13.3 & 14.0 & 12.0 & 14.0\\
  19 & 14.3 & 14.5 & 14.8 & 12.3 & 13.0 & 14.0 & 12.2 & 13.6\\
  20 & 13.5 & 13.8 & 14.3 & 12.0 & 13.0 & 14.2 & 12.7 & 13.4\\ \midrule
  21 & 12.8 & 13.3 & 14.0 & 12.0 & 13.2 & 14.7 & 13.4 & 13.3\\
  22 & 12.3 & 13.0 & 14.0 & 12.2 & 13.7 & 15.4 & 14.2 & 13.6\\
  23 & 12.0 & 13.0 & 14.2 & 12.7 & 14.4 & 16.2 & 15.2 & 14.0\\
  24 & 12.0 & 13.2 & 14.7 & 13.4 & 15.2 & 17.2 & 16.3 & 14.6\\
  25 & 12.2 & 13.7 & 15.4 & 14.2 & 16.2 & 18.3 & 17.5 & 15.4\\
  26 & 12.7 & 14.4 & 16.2 & 15.2 & 17.3 & 19.5 & 18.7 & 16.3\\
  27 & 13.4 & 15.2 & 17.2 & 16.3 & 18.5 & 20.7 & 19.9 & 17.3\\
  28 & 14.2 & 16.2 & 18.3 & 17.5 & 19.7 & 21.9 & 21.0 & 18.4\\
  29 & 15.2 & 17.3 & 19.5 & 18.7 & 20.9 & 23.0 & 21.9 & 19.5\\
  30 & 16.3 & 18.5 & 20.7 & 19.9 & 22.0 & 23.9 & 22.8 & 20.6\\ \midrule
  31 & 17.5 & 19.7 & 21.9 & 21.0 & 22.9 & 24.8 & 23.4 & 21.6\\
  32 & 18.7 & 20.9 & 23.0 & 21.9 & 23.8 & 25.4 & 23.8 & 22.5\\
  33 & 19.9 & 22.0 & 23.9 & 22.8 & 24.4 & 25.8 & 24.0 & 23.2\\
  34 & 21.0 & 22.9 & 24.8 & 23.4 & 24.8 & 26.0 & 23.9 & 23.8\\
  35 & 21.9 & 23.8 & 25.4 & 23.8 & 25.0 & 25.9 & 23.6 & 24.2\\
  36 & 22.8 & 24.4 & 25.8 & 24.0 & 24.9 & 25.6 & 23.1 & 24.4\\
  37 & 23.4 & 24.8 & 26.0 & 23.9 & 24.6 & 25.1 & 22.4 & 24.3\\
  38 & 23.8 & 25.0 & 25.9 & 23.6 & 24.1 & 24.4 & 21.5 & 24.1\\
  39 & 24.0 & 24.9 & 25.6 & 23.1 & 23.4 & 23.5 & 20.5 & 23.6\\
  40 & 23.9 & 24.6 & 25.1 & 22.4 & 22.5 & 22.5 & 19.3 & 22.9\\ \midrule
  41 & 23.6 & 24.1 & 24.4 & 21.5 & 21.5 & 21.3 & 18.1 & 22.1\\
  42 & 23.1 & 23.4 & 23.5 & 20.5 & 20.3 & 20.1 & 17.0 & 21.1\\
  43 & 22.4 & 22.5 & 22.5 & 19.3 & 19.1 & 19.0 & 15.8 & 20.1\\
  44 & 21.5 & 21.5 & 21.3 & 18.1 & 18.0 & 17.8 & 14.7 & 19.0\\
  45 & 20.5 & 20.3 & 20.1 & 17.0 & 16.8 & 16.7 & 13.8 & 17.9\\
  46 & 19.3 & 19.1 & 19.0 & 15.8 & 15.7 & 15.8 & 13.0 & 16.8\\
  47 & 18.1 & 18.0 & 17.8 & 14.7 & 14.8 & 15.0 & 12.5 & 15.8\\
  48 & 17.0 & 16.8 & 16.7 & 13.8 & 14.0 & 14.5 & 12.1 & 15.0\\
  49 & 15.8 & 15.7 & 15.8 & 13.0 & 13.5 & 14.1 & 12.0 & 14.3\\
  50 & 14.7 & 14.8 & 15.0 & 12.5 & 13.1 & 14.0 & 12.1 & 13.8\\ \midrule
  51 & 13.8 & 14.0 & 14.5 & 12.1 & 13.0 & 14.1 & 12.5 & 13.4\\
  52 & 13.0 & 13.5 & 14.1 & 12.0 & 13.1 & 14.5 & 13.1 & 13.3\\
  53 & 12.5 & 13.1 & 14.0 & 12.1 & 13.5 & 15.1 & 13.8 & 13.4\\
  54 & 12.1 & 13.0 & 14.1 & 12.5 & 14.1 & 15.8 & 14.8 & 13.8\\
  55 & 12.0 & 13.1 & 14.5 & 13.1 & 14.8 & 16.8 & 15.9 & 14.3\\
  56 & 12.1 & 13.5 & 15.1 & 13.8 & 15.8 & 17.9 & 17.0 & 15.0\\
  57 & 12.5 & 14.1 & 15.8 & 14.8 & 16.9 & 19.0 & 18.2 & 15.9\\
  58 & 13.1 & 14.8 & 16.8 & 15.9 & 18.0 & 20.2 & 19.4 & 16.9\\
  59 & 13.8 & 15.8 & 17.9 & 17.0 & 19.2 & 21.4 & 20.5 & 17.9\\
  60 & 14.8 & 16.9 & 19.0 & 18.2 & 20.4 & 22.5 & 21.6 & 19.0\\ \midrule
\end{longtable}
\end{landscape}

\section*{7. Tabla girada con notas}
\begin{sidewaystable}
\centering
\begin{threeparttable}
\caption{Tabla girada 90° con \texttt{rotating} y notas.}
\begin{tabular}{lrrrrrr}
  \toprule
  & \multicolumn{3}{c}{Grupo A\tnote{1}} & \multicolumn{3}{c}{Grupo B\tnote{2}}\\
  \cmidrule(lr){2-4}\cmidrule(lr){5-7}
  Variable & Media & DE & $n$ & Media & DE & $n$\\
  \midrule
  Edad & 21,4 & 2,1 & 30 & 22,0 & 1,8 & 28\\
  Nota final & 13,2 & 3,4 & 30 & 11,9 & 4,0 & 28\\
  Horas de estudio & 9,5 & 3,0 & 30 & 7,1 & 2,6 & 28\\
  \bottomrule
\end{tabular}
\begin{tablenotes}\small
  \item[1] Turno mañana. \item[2] Turno noche.
\end{tablenotes}
\end{threeparttable}
\end{sidewaystable}

\section*{8. Tabla de 10 columnas que no cabe (se reduce sola)}
\noindent\resizebox{\linewidth}{!}{%
\begin{tabular}{*{10}{c}}
  \toprule
  $x_1$ & $x_2$ & $x_3$ & $x_4$ & $x_5$ & $x_6$ & $x_7$ & $x_8$ & $x_9$ & $x_{10}$\\
  \midrule
  $\frac{1}{2}$ & $\sqrt{2}$ & $e^{i\pi}$ & $\int_0^1 x\,dx$ & $\sum_{k=1}^{n} k$ & $\lim_{x\to0}\frac{\sin x}{x}$ & $\binom{n}{k}$ & $\nabla\cdot\vec E$ & $\oint_C \vec B\cdot d\vec l$ & $\aleph_0$\\
  \bottomrule
\end{tabular}}
\end{document}
